\begin{abstract}
    As doenças cardiovasculares são a principal causa de morte no mundo, e a cardiopatia congênita é a mais comum entre as doenças cardíacas em pacientes pediátricos. Técnicas de inteligência artificial têm sido amplamente utilizadas para auxiliar no diagnóstico de doenças cardíacas, principalmente devido à sua capacidade de processar grandes quantidades de dados e identificar padrões. Sabendo disso, este trabalho avalia o desempenho de métodos de aprendizado de máquina e redes neurais artificias na detecção de patologias cardíacas em pacientes pediátricos. Para tal, foi utilizado um conjunto de dados com informações clínicas de pacientes pediátricos com e sem patologias cardíacas. Os resultados obtidos demonstram que modelos de aprendizado de máquina e redes neurais artificiais são capazes de identificar patologias cardíacas em pacientes pediátricos com alta precisão, especialmente quando as etapas de analise exploratória e pré-processamento dos dados são adequadamente realizados.
\end{abstract}

\begin{IEEEkeywords}
    doenças cardiovasculares; pediatria; inteligência artificial; classificação;

\end{IEEEkeywords}